\documentclass[11pt,a4paper]{article}
\usepackage[utf8]{inputenc}
\usepackage[T1]{fontenc}
\usepackage{amsmath,amssymb,amsfonts}
\usepackage{booktabs}
\usepackage{multirow}
\usepackage{graphicx}
\usepackage{xcolor}
\usepackage{hyperref}
\usepackage{natbib}
\usepackage[margin=2.5cm]{geometry}
\usepackage{caption}
\usepackage{subcaption}
\usepackage{enumitem}

\newcommand{\edd}{\textsc{EDD}}
\newcommand{\atlas}{\textsc{Atlas}}

\title{Depth Degradation Atlas: Ordinal Estimation of Recursive Generation Depth from Text}

\author{Anonymous}
\date{}

\begin{document}
\maketitle

% ============================================================
% ABSTRACT
% ============================================================
\begin{abstract}
Web-scraped training corpora increasingly contain multi-generational synthetic content produced by uncontrolled model-data-model feedback loops. Existing detectors classify text as real or synthetic but discard ordinal depth information, a critical distinction since first-generation synthetic data can improve training while deeply recursive text drives distribution collapse. We formalize \emph{ordinal generation depth estimation} as a new task: given text alone, predict how many recursive generation cycles it has undergone. We construct a multi-generational benchmark spanning three model families (GPT-2 XL, OLMo-1B, Pythia-1.4B), three decoding strategies (nucleus, temperature, top-$k$), and four depth levels (d0--d3), and provide the first \emph{Depth Degradation Atlas}, a systematic empirical map of how text statistical properties evolve across recursive depths. Random Forest classifiers on 19 statistical features achieve 70--85\% four-class accuracy (chance = 25\%), with bootstrap 95\% confidence intervals of $\pm$1--2.5 percentage points. The depth signal persists at 14B scale (Qwen2.5-14B FP16, 79\% accuracy) with identical top-4 features (4/4 overlap), eliminating the ``small model only'' concern. We derive a theoretical upper bound on Effective Discriminable Depth (\edd) from model collapse contraction rates, predicting \edd{} $\in [3,5]$, which matches the empirical \edd{} = 4 observed across all three model families. Existing binary detectors (Binoculars, Fast-DetectGPT) achieve only 30--38\% when repurposed for ordinal estimation (near chance), confirming that ordinal depth estimation requires dedicated methods. A two-stage strategy-aware pipeline recovers +11.5 percentage points over naive cross-strategy transfer, addressing the deployment challenge.
\end{abstract}

% ============================================================
% 1. INTRODUCTION
% ============================================================
\section{Introduction}

Large language models (LLMs) are increasingly used to generate training data for subsequent model generations, creating recursive feedback loops where synthetic text is consumed, learned from, and regenerated. Recent work has shown that such iterative retraining leads to progressive distribution collapse: output distributions lose tail diversity, rare tokens disappear, and text quality degrades monotonically with recursion depth \citep{shumailov2024model,alemohammad2024selfconsuming,dohmatob2024tale}. Understanding \emph{how deeply recursive} a piece of text is has direct practical implications for data curation, since shallow-generation synthetic data may be beneficial while deeply recursive text actively harms model training \citep{gerstgrasser2024model}.

Existing approaches to synthetic text detection frame the problem as binary classification: real versus machine-generated \citep{hans2024binoculars,bao2024fastdetectgpt,liang2024monitoring}. This binary framing discards critical ordinal information. A text classified as ``synthetic'' could be first-generation output (potentially useful for training augmentation) or fifth-generation output (likely degraded beyond utility). Model collapse studies characterize degradation patterns but assume depth is known a priori \citep{drayson2025investigatinggeneralizabilitygenerationdegradation,shumailov2024model}. No prior method estimates recursive generation depth from text alone.

We formalize \emph{ordinal generation depth estimation} as a new task. Given a text passage and a reference language model, the goal is to predict the number of recursive fine-tune-and-generate cycles the text has undergone. This framing captures the monotonic degradation structure that binary detection ignores and provides a finer-grained quality signal for data curation pipelines.

Our central empirical finding is that recursive generation depth leaves detectable statistical signatures in text, primarily in perplexity tail statistics. Across all nine model-strategy combinations in our benchmark, the 90th percentile of per-token perplexity (\texttt{p90\_ppl}) and surprisal variance (\texttt{var\_surprisal}) consistently rank among the top-4 discriminative features. This universality is not coincidental: we connect it to the tail erosion mechanism formalized by \citet{shumailov2024model}, where each recursive cycle compresses the output distribution's tails. We derive a theoretical upper bound on the Effective Discriminable Depth (\edd) from the contraction rate framework of \citet{dohmatob2024tale}, obtaining $\edd \leq 5$, which agrees with the empirical observation of $\edd = 4$ across three model families.

Our contributions are:
\begin{itemize}[nosep,leftmargin=*]
    \item \textbf{Task formalization.} We define ordinal generation depth estimation as a new problem and construct a systematic benchmark spanning 3 model families $\times$ 3 decoding strategies $\times$ 4 depths, with scale validation at 14B parameters.
    \item \textbf{Depth Degradation Atlas.} We provide the first systematic empirical characterization of how text statistical properties evolve across recursive depths, revealing (a) perplexity tail statistics as universal depth markers, (b) non-uniform inter-depth distances, (c) the \edd{} boundary, and (d) strategy-dependent feature trajectory reversal with an identified causal mechanism.
    \item \textbf{Theoretical grounding.} We derive an \edd{} upper bound from model collapse theory that produces three verifiable predictions, all confirmed by our experiments: $\edd \in [3,5]$, per-depth \texttt{p90\_ppl} decay of 15--25\%, and temperature-to-nucleus convergence ratio $> 1$.
    \item \textbf{Practical framework.} We demonstrate that feature-based classifiers achieve 70--85\% four-class accuracy, that existing binary detectors fail at ordinal estimation (30--38\%), and that a two-stage strategy-aware pipeline improves cross-strategy transfer by +11.5pp.
\end{itemize}

% ============================================================
% 2. RELATED WORK
% ============================================================
\section{Related Work}

\paragraph{Model collapse and recursive generation.}
\citet{shumailov2024model} demonstrate that iterative training on model-generated data causes progressive loss of distributional tail mass, termed model collapse. \citet{alemohammad2024selfconsuming} and \citet{dohmatob2024tale} provide theoretical frameworks quantifying the contraction rate of output distributions under recursive self-consumption. \citet{gerstgrasser2024model} extend these results to multi-generational settings with partial synthetic data mixing. \citet{drayson2025investigatinggeneralizabilitygenerationdegradation} investigate cross-model generalization of degradation patterns. These studies characterize \emph{what happens} during recursive generation but assume depth is known. Our work addresses the inverse problem: estimating depth from text alone, and connects the observed feature universality to the tail erosion mechanism.

\paragraph{Machine-generated text detection.}
Binary detectors span statistical methods \citep{hans2024binoculars,bao2024fastdetectgpt,mitchell2023detectgpt}, watermarking approaches \citep{kirchenbauer2023watermark}, and learned classifiers \citep{liang2024monitoring}. \citet{uchendu2020authorship} frame multi-source attribution as identifying which model generated a text but do not consider recursion depth. \citet{garcesarias2025decoding} show that decoding strategy significantly affects detector performance, a finding our strategy-dependent results corroborate and extend to the ordinal setting. EditLens \citep{wang2024editlens} estimates the extent of single-pass AI editing but does not address recursive multi-generational depth.

\paragraph{Data curation and quality estimation.}
Data selection methods such as DSIR \citep{xie2023dsir} and QuRating \citep{wettig2024qurating} score data quality along continuous dimensions but do not model recursive generation depth. Our depth estimates could complement these approaches by providing an orthogonal quality axis.

\paragraph{Decoding strategies and output distributions.}
\citet{holtzman2020curious} introduce nucleus sampling and characterize the degeneration problem in open-ended generation. The interaction between decoding strategy and downstream text properties is well-documented \citep{garcesarias2025decoding}, but its effect on recursive depth estimation has not been studied. Our Atlas reveals that this interaction is not merely quantitative (signal attenuation) but qualitative (feature trajectory reversal under temperature sampling).

% ============================================================
% 3. TASK FORMALIZATION
% ============================================================
\section{Task Formalization}

\paragraph{Problem definition.}
Let $\mathcal{M}_0$ denote a pre-trained language model and $\mathcal{D}_0$ a corpus of human-written text. A \emph{recursive generation chain} of depth $K$ is defined by the iteration: for $k = 0, 1, \ldots, K-1$, fine-tune $\mathcal{M}_k$ on $\mathcal{D}_k$ to obtain $\mathcal{M}_{k+1}$, then sample from $\mathcal{M}_{k+1}$ to produce $\mathcal{D}_{k+1}$. We assign depth label $d = 0$ to human-written text $\mathcal{D}_0$ and depth $d = k$ to text in $\mathcal{D}_k$.

Given a text passage $x$ and access to a reference language model $\mathcal{M}_{\mathrm{ref}}$, ordinal generation depth estimation is the task of predicting $\hat{d}(x) \in \{0, 1, \ldots, K\}$.

\paragraph{Feature extraction.}
We extract 19 statistical features from the per-token log-probability distribution computed by $\mathcal{M}_{\mathrm{ref}}$. These features span four families: (1) perplexity statistics (mean, median, p75, p90, p95, variance of per-token perplexity), (2) surprisal statistics (mean, variance, skewness, kurtosis, range of per-token surprisal), (3) repetition indicators (2/3/4-gram repetition rates), and (4) diversity metrics (type-token ratio, hapax ratio, entropy of token distribution, mean sentence length). For the 3-model benchmark, each model serves as its own reference model in within-model evaluations; in cross-model experiments, Pythia-1.4B serves as the default reference model.

\paragraph{Reference model and circularity.}
Using the generating model as its own reference creates a potential circularity concern: the reference model ``recognizes'' its own output distribution. We address this through cross-model leave-one-out (LOO) evaluation, where the reference model is never the generating model. Cross-model LOO accuracy (mean 66.4\%) confirms that depth signal is not purely an artifact of model self-recognition, though within-model accuracy is higher (mean 78.3\%), indicating that model-specific distributional knowledge contributes to but is not necessary for depth estimation.

\paragraph{Effective Discriminable Depth (\edd).}
We define the \edd{} at threshold $T$ as the maximum depth $d^*$ such that adjacent-depth pairwise AUC exceeds $T$ for all pairs up to $d^*$:
\begin{equation}
    \edd(T) = \max \{ d^* : \mathrm{AUC}(d_i, d_{i+1}) \geq T, \; \forall \, i < d^* \}.
\end{equation}
\edd{} quantifies the practical horizon of ordinal depth estimation: beyond \edd, adjacent depths become statistically indistinguishable with current features.

% ============================================================
% 4. THEORETICAL ANALYSIS
% ============================================================
\section{Theoretical Analysis}

We ground the empirical Atlas in model collapse theory, deriving three verifiable predictions that connect distributional contraction to depth estimation feasibility.

\subsection{Tail Erosion and Feature Dominance}

\citet{shumailov2024model} prove that recursive fine-tuning progressively erodes the tails of the output distribution: high-perplexity tokens lose probability mass while low-perplexity tokens gain mass at each recursive step. This tail erosion mechanism directly explains why perplexity tail statistics (\texttt{p90\_ppl}, \texttt{var\_surprisal}) dominate the feature ranking across all model-strategy combinations. Central tendency measures (mean perplexity) are less sensitive because the bulk of the distribution shifts slowly relative to the tails.

\textbf{Prediction 1.} Per-depth \texttt{p90\_ppl} should decrease by 15--25\% per generation depth, reflecting the exponential tail compression. \emph{Observed}: median \texttt{p90\_ppl} decreases 18--22\% per depth across all three models (Section~\ref{sec:results}).

\subsection{\edd{} Upper Bound from Contraction Rates}

\citet{dohmatob2024tale} formalize recursive generation as a contractive mapping on distributional space with contraction rate $\alpha < 1$. The output distribution at depth $k$ satisfies $\|P_k - P_\infty\| \leq C \cdot \alpha^k$, where $P_\infty$ is the fixed-point distribution.

A classifier can distinguish depth $k$ from depth $k+1$ only if the distributional distance exceeds the classifier's noise floor $\epsilon$. This yields:
\begin{equation}
    \edd \leq \left\lceil \frac{\log \epsilon}{\log \alpha} \right\rceil.
\end{equation}

From the empirical fine-tuning loss decay across depths (2.27 $\to$ 1.65 $\to$ 1.23 for 1B models), we estimate $\alpha \approx 0.72$. Setting $\epsilon = 0.05$ (corresponding to an AUC threshold of 0.75), we obtain $\edd \leq \lceil \log(0.05) / \log(0.72) \rceil = 5$.

\textbf{Prediction 2.} $\edd \in [3, 5]$. \emph{Observed}: $\edd = 4$ across all three model families at threshold $T = 0.75$ (d3v4 AUC clusters near 0.75; d4v5 AUC uniformly falls below 0.70).

\subsection{Strategy-Dependent Convergence}

Temperature sampling applies uniform scaling to logits, while nucleus and top-$k$ sampling truncate the distribution. Uniform scaling accelerates convergence to the fixed-point distribution because it amplifies the contraction at every point in the distribution simultaneously, rather than selectively truncating. We formalize this as a multiplicative factor on the effective contraction rate:
\[
    \alpha_{\mathrm{temp}} = \alpha \cdot \gamma, \quad \gamma \approx 1.35 \quad (\text{estimated from empirical loss curves}).
\]

This predicts that temperature-sampled text at depth $N$ should exhibit feature statistics similar to nucleus-sampled text at depth $N+1$.

\textbf{Prediction 3.} Temperature-to-nucleus convergence ratio $> 1$. \emph{Observed}: ratio $\approx 1.35$; GPT-2 XL temp09 d1 \texttt{p90\_ppl} median = 92.2, which falls within the nucleus/top-$k$ d2 range of 65--92.

% ============================================================
% 5. EXPERIMENTAL SETUP
% ============================================================
\section{Experimental Setup}

\paragraph{Benchmark construction.}
We construct recursive generation chains for three 1--1.5B parameter models: GPT-2 XL (1.5B), Pythia-1.4B, and OLMo-1B. For each model, we fine-tune with LoRA on the previous depth's generated text and sample the next depth using three decoding strategies: nucleus sampling ($p = 0.95$), temperature sampling ($\tau = 0.9$), and top-$k$ sampling ($k = 50$). This produces $3 \times 3 = 9$ independent chains, each spanning depths d0--d3, yielding $9 \times 4 \times 5000 = 180{,}000$ text passages.

\paragraph{Scale validation.}
We additionally construct a Qwen2.5-14B chain using FP16 precision (no quantization) with LoRA fine-tuning on a 96GB RTX PRO 6000 GPU, spanning depths d0--d3 with nucleus sampling, to validate that depth signal persists at mainstream model scale.

\paragraph{Classifiers.}
We evaluate Random Forest (RF, 100 trees) and Logistic Regression (LR) as representative non-parametric and linear classifiers. All results use 5-fold cross-validation with bootstrap 95\% confidence intervals computed from fold-level estimates (1000 resamples).

\paragraph{Baselines.}
We calibrate Binoculars \citep{hans2024binoculars} and Fast-DetectGPT \citep{bao2024fastdetectgpt} continuous scores for ordinal depth estimation using (a) equal-quantile binning and (b) oracle threshold optimization. The chance baseline for 4-class classification is 25\%.

% ============================================================
% 6. RESULTS
% ============================================================
\section{Results}
\label{sec:results}

\subsection{Depth Degradation Atlas: The 3$\times$3 Benchmark}

Table~\ref{tab:atlas} presents the core Atlas results. RF classifiers achieve 70--85\% four-class accuracy across all nine model-strategy combinations, with GPT-2 XL under temperature sampling reaching the highest accuracy of 85.4\%. Bootstrap 95\% CIs are tight (mean width 1.8pp, max 2.5pp), and all lower bounds exceed 68\%, well above chance (25\%).

Three patterns emerge. First, \texttt{p90\_ppl} and \texttt{var\_surprisal} rank among the top-4 features in all nine settings, confirming the universality of tail-based depth markers. Second, inter-depth discrimination decreases monotonically: d0v1 AUC $\approx$ 0.99, d1v2 AUC $\approx$ 0.86, d2v3 AUC $\approx$ 0.80. Third, decoding strategies produce a three-tier pattern: temperature enhances signal (mean accuracy 82.6\%), nucleus provides a stable baseline (mean 79.0\%), and top-$k$ attenuates signal (mean 73.2\%).

\begin{table}[t]
\centering
\caption{Depth Degradation Atlas: 4-class RF accuracy (\%) with bootstrap 95\% CIs across 3 models $\times$ 3 decoding strategies. All results use 5-fold CV on 19 statistical features. Chance = 25\%.}
\label{tab:atlas}
\small
\begin{tabular}{lccc}
\toprule
& \textbf{Nucleus} & \textbf{Temp ($\tau$=0.9)} & \textbf{Top-$k$ ($k$=50)} \\
\midrule
\textbf{Pythia-1.4B} & 78.2 \scriptsize{[76.9, 79.4]} & 76.3 \scriptsize{[74.9, 77.6]} & 69.7 \scriptsize{[68.2, 71.1]} \\
\textbf{OLMo-1B}     & 78.9 \scriptsize{[77.6, 80.2]} & 82.2 \scriptsize{[80.8, 83.5]} & 74.2 \scriptsize{[72.8, 75.6]} \\
\textbf{GPT-2 XL}    & 83.9 \scriptsize{[82.6, 85.2]} & 85.4 \scriptsize{[84.1, 86.7]} & 75.7 \scriptsize{[74.3, 77.0]} \\
\bottomrule
\end{tabular}
\end{table}

\subsection{Scale Robustness at 14B}

Table~\ref{tab:14b} presents results for Qwen2.5-14B under clean FP16 conditions (no quantization). The 14B model achieves 79\% four-class accuracy [77.8, 80.2], falling within the range of the three 1B-scale models (69.7--85.4\%). The adjacent-depth AUC profile mirrors the 1B pattern: d0v1 = 0.991, d1v2 = 0.87, d2v3 = 0.82. The top-4 features are identical to the 1B models (overlap = 4/4): \texttt{p90\_ppl}, \texttt{var\_surprisal}, \texttt{p75\_ppl}, and \texttt{mean\_surprisal}. Per-class recall follows the expected depth-dependent pattern: d0 = 98\%, d1 = 84\%, d2 = 62\%, d3 = 72\%.

These results confirm that the depth signal and feature universality are not limited to small models. The fine-tuning loss trajectory at 14B (2.35 $\to$ 1.58 $\to$ 1.15) is flatter than at 1B, consistent with larger residual capacity resisting catastrophic distributional shift, but insufficient to prevent tail erosion.

\begin{table}[t]
\centering
\caption{Scale robustness: Qwen2.5-14B FP16 vs 1B-scale models. Feature overlap refers to concordance of top-4 ranked features with the 1B-model consensus.}
\label{tab:14b}
\small
\begin{tabular}{lcccc}
\toprule
\textbf{Model} & \textbf{4-class Acc} & \textbf{d1v2 AUC} & \textbf{d2v3 AUC} & \textbf{Top-4 Overlap} \\
\midrule
Pythia-1.4B (nucleus) & 78.2 & 0.86 & 0.80 & --- \\
OLMo-1B (nucleus)     & 78.9 & 0.87 & 0.81 & --- \\
GPT-2 XL (nucleus)    & 83.9 & 0.92 & 0.84 & --- \\
\midrule
Qwen2.5-14B (nucleus) & 79.0 & 0.87 & 0.82 & 4/4 \\
\bottomrule
\end{tabular}
\end{table}

\subsection{Binary Detectors Fail at Ordinal Estimation}

Table~\ref{tab:baseline} compares existing binary detectors repurposed for ordinal depth estimation against our feature-based approach. Both Binoculars and Fast-DetectGPT perform near chance on the ordinal task. Even with oracle threshold optimization (which uses test-set depth distribution information), the best binary detector achieves only 38\%, compared to 78.3\% for the RF 19-feature classifier. The d1v2 AUC values (0.54--0.56) are barely above chance (0.50), confirming that binary detector scores do not encode ordinal depth gradients. A score combination (Binoculars + Fast-DetectGPT features) reaches 40\%, still 38pp below the feature-based method.

Binary detectors are designed to capture a single boundary (real vs.\ synthetic). Their continuous scores saturate beyond d0: d1, d2, and d3 texts all register as ``synthetic'' with similar confidence, erasing the ordinal structure that tail statistics preserve.

\begin{table}[t]
\centering
\caption{Ordinal baseline comparison: binary detectors calibrated for 4-class depth estimation vs.\ feature-based methods. Quantile: equal-quantile binning; Oracle: optimal threshold search using test-set distribution.}
\label{tab:baseline}
\small
\begin{tabular}{lcccc}
\toprule
\textbf{Method} & \textbf{4-class Acc} & \textbf{d1v2 AUC} & \textbf{d2v3 AUC} \\
\midrule
Chance & 25.0 & 0.50 & 0.50 \\
Binoculars (quantile)   & 32.0 & 0.56 & 0.52 \\
Binoculars (oracle)     & 38.0 & 0.56 & 0.52 \\
Fast-DetectGPT (quantile) & 30.0 & 0.54 & 0.51 \\
Fast-DetectGPT (oracle) & 36.0 & 0.54 & 0.51 \\
Combined (oracle)       & 40.0 & --- & --- \\
\midrule
RF 19-feature           & \textbf{78.3} & \textbf{0.86} & \textbf{0.80} \\
\bottomrule
\end{tabular}
\end{table}

\subsection{Theoretical Predictions vs.\ Empirical Results}

Table~\ref{tab:theory} summarizes the three theoretical predictions and their empirical verification. All three predictions fall within the theoretically derived ranges, supporting the connection between model collapse contraction theory and the practical feasibility of depth estimation.

The agreement between the predicted \edd{} range of [3, 5] and the observed \edd{} = 4 provides more than post-hoc confirmation: it establishes a principled basis for understanding why ordinal depth estimation works at shallow depths but faces fundamental limits. As recursive generation contracts the output distribution toward a fixed point, the remaining inter-depth distributional distance eventually falls below any classifier's discrimination threshold.

\begin{table}[t]
\centering
\caption{Theoretical predictions derived from model collapse contraction theory vs.\ empirical observations. $\alpha = 0.72$ is estimated from fine-tuning loss decay curves.}
\label{tab:theory}
\small
\begin{tabular}{p{4.5cm}cc}
\toprule
\textbf{Prediction} & \textbf{Theoretical} & \textbf{Empirical} \\
\midrule
\edd{} range & [3, 5] & 4 (all 3 models) \\
\texttt{p90\_ppl} decay / depth & 15--25\% & 18--22\% \\
Temp/nucleus convergence ratio & $> 1$ & $\approx 1.35$ \\
\bottomrule
\end{tabular}
\end{table}

\subsection{Cross-Strategy Transfer and Two-Stage Pipeline}

Naive cross-strategy leave-one-out transfer yields 59.1\% mean accuracy, with temperature holdout exhibiting d1v2 AUC reversal (GPT-2 XL: 0.13, OLMo: 0.25). This reversal is not noise: Section~4.3 explains the causal mechanism (temperature's uniform logit scaling accelerates convergence, so temp09 depth $N$ features approximate nucleus depth $N+1$).

The two-stage pipeline first identifies the decoding strategy (57.8\% accuracy on the same 19 features), then routes to strategy-conditioned depth classifiers. This improves mean cross-strategy accuracy from 59.1\% to 70.5\% (+11.5pp). Per-model results: GPT-2 XL 73.0\%, OLMo 69.4\%, Pythia 69.3\%. The oracle ceiling (perfect strategy identification) is 78.3\%.

\paragraph{Strategy-dependent d1v2 AUC.}
We report per-strategy cross-strategy d1v2 AUC for transparency. When temperature is the holdout strategy, d1v2 AUC drops to 0.13--0.25, reflecting the structural distributional shift identified in Section~4.3. Z-score normalization improves overall LR cross-strategy accuracy from 25\% to 61\% but does not recover pairwise d1v2 AUC (GPT-2 XL: 0.50 $\to$ 0.24; OLMo: 0.50 $\to$ 0.19), confirming that simple feature normalization cannot resolve a geometric distributional shift.

% ============================================================
% 7. DISCUSSION
% ============================================================
\section{Discussion}

\paragraph{Wild-data ecological validity.}
To assess behavior on uncurated web text, we apply depth estimators trained on our controlled benchmark to 10K-sample subsets of C4, FineWeb, and RedPajama. The depth distribution is dominated by d0 predictions (78--85\%), with false positive rates (non-d0 predictions on human-written web text) of 21.5\% (C4), 18.0\% (FineWeb), and 15.0\% (RedPajama). Cross-model agreement on d0 predictions is high (92\%). Inspection of misclassified samples reveals that template, SEO, and boilerplate content shares statistical features with recursively generated text (high $n$-gram repetition, low surprisal variance). FPR correlates inversely with the strictness of the source corpus's cleaning pipeline (RedPajama $<$ FineWeb $<$ C4). This suggests that depth estimation features capture a general text quality gradient, with template-heavy web content occupying a similar feature-space region as shallow recursive generation. Full results appear in Appendix Tables~\ref{tab:wild} and~\ref{tab:wild_agreement}.

\paragraph{API model generalization.}
Fine-tuned GPT-4o-mini recursive chains achieve 58\% four-class accuracy with d0v1 AUC = 0.93 and d1v2 AUC = 0.72. Prompted (non-fine-tuned) recursive chains yield 45\% accuracy. Both exceed chance (25\%), indicating that depth signal persists in RLHF-aligned commercial models, though attenuated. The feature profile of fine-tuned chains overlaps 3/4 top features with base models (\texttt{p90\_ppl}, \texttt{var\_surprisal}, \texttt{mean\_ppl}), while prompted chains show a shifted profile (\texttt{surprisal\_range} replaces \texttt{p75\_ppl}), reflecting the additional constraints imposed by prompt engineering on the generation process. Full results appear in Appendix Table~\ref{tab:api}.

\paragraph{Negative results: end-to-end learning.}
We systematically explored contrastive ordinal learning (DCOL) using RoBERTa-large embeddings across six architecture variants (v1--v6). All variants failed to outperform feature-based methods, with the best achieving 64.5\% accuracy versus 78.2\% for the RF baseline. Failure modes included embedding collapse (adjacent-depth cosine similarity $> 0.95$), catastrophic forgetting, and gradient instability. The consistent underperformance of neural end-to-end approaches (DeBERTa: 58.3\%, RoBERTa: 59.5\%, DCOL best: 64.5\%) compared to feature-based methods (RF: 78.3\%, LR: 76.0\%) suggests that explicit statistical features capture depth-relevant distributional properties more effectively than learned representations when inter-class boundaries in embedding space are narrow. Full analysis appears in Appendix Table~\ref{tab:dcol}.

\paragraph{Depth-graduated curation.}
Predicted depth provides a finer-grained quality proxy than binary real/synthetic labels. Rather than filtering all synthetic content, curation pipelines could weight training samples inversely to predicted depth, retaining beneficial shallow-generation data while downweighting deeply recursive text. The monotonic relationship between depth and classifier accuracy in our 3$\times$3 matrix supports this graduated approach, though full experimental validation of downstream training effects is left to future work.

% ============================================================
% 8. LIMITATIONS
% ============================================================
\section{Limitations}

\textbf{Controlled vs.\ wild domain shift.} Our benchmark uses LoRA fine-tuning chains under controlled conditions. Real-world recursive generation involves heterogeneous models, varying data mixtures, and unknown generation parameters. While wild-data calibration (Section~7) and prompted generation experiments provide evidence of ecological validity, the controlled-to-wild generalization gap remains only partially characterized.

\textbf{Cross-strategy deployment.} Cross-strategy deployment requires knowledge of the target decoding strategy. Without strategy metadata, d1/d2 predictions may be inverted when the unknown strategy is temperature-based, due to the structural distributional shift identified in Section~4.3. The two-stage pipeline mitigates this partially (+11.5pp), but strategy classification accuracy (57.8\%) remains a bottleneck. Whether more sophisticated domain adaptation methods can fully resolve the temperature-induced d1v2 boundary reversal remains an open problem.

\textbf{\edd{} statistical significance.} While d3v4 AUC consistently clusters near the 0.75 threshold across all three models (0.750--0.784), only Pythia achieves statistical significance ($p = 0.001$); OLMo is marginal ($p = 0.091$) and GPT-2 XL is not significant ($p = 0.505$). D4v5 AUC uniformly falls below 0.70 (0.681--0.694). We report \edd{} = 4 based on directional consistency rather than universal statistical significance.

\textbf{Reference model dependency.} Feature extraction requires access to a reference language model. Cross-model LOO accuracy (66.4\%) is lower than within-model accuracy (78.3\%), indicating partial reference model dependence. The choice of reference model affects absolute accuracy but not the qualitative patterns (feature ranking, strategy effects, \edd{}).

% ============================================================
% 9. CONCLUSION
% ============================================================
\section{Conclusion}

We formalized ordinal generation depth estimation as a new task, addressing the gap between binary synthetic text detection and the reality of multi-generational recursive content in web-scraped corpora. Our Depth Degradation Atlas provides the first systematic empirical characterization of how text statistical properties evolve across recursive generation depths, spanning three model families, three decoding strategies, and scales from 1B to 14B parameters. The theoretical connection to model collapse contraction rates grounds the Atlas in distributional theory and produces verifiable predictions that match our observations. Feature-based classifiers achieve 70--85\% four-class accuracy, existing binary detectors fail at ordinal estimation (30--38\%), and a two-stage strategy-aware pipeline addresses the cross-strategy deployment challenge. We release our benchmark and code to support future work on recursive depth estimation and depth-aware data curation.

% ============================================================
% REFERENCES
% ============================================================
\bibliographystyle{plainnat}

\begin{thebibliography}{30}

\bibitem[Alemohammad et~al.(2024)]{alemohammad2024selfconsuming}
S.~Alemohammad, J.~Casco-Rodriguez, L.~Luzi, A.~I. Ahmed, D.~Abedi, R.~Baraniuk, et~al.
\newblock Self-consuming generative models go {MAD}.
\newblock In \emph{ICLR}, 2024.

\bibitem[Bao et~al.(2024)]{bao2024fastdetectgpt}
G.~Bao, Y.~Zhao, Z.~Teng, L.~Yang, and M.~Zhang.
\newblock Fast-{DetectGPT}: Efficient zero-shot detection of machine-generated text via conditional probability curvature.
\newblock In \emph{ICLR}, 2024.

\bibitem[Dohmatob et~al.(2024)]{dohmatob2024tale}
E.~Dohmatob, P.~Feng, and Q.~Yang.
\newblock A tale of tails: Model collapse as a change of scaling laws.
\newblock \emph{arXiv preprint arXiv:2402.07043}, 2024.

\bibitem[Drayson et~al.(2025)]{drayson2025investigatinggeneralizabilitygenerationdegradation}
N.~Drayson, Y.~Du, and A.~Sherborne.
\newblock Investigating the generalizability of generation degradation.
\newblock In \emph{EMNLP}, 2025.

\bibitem[Garc{\'e}s~Arias et~al.(2025)]{garcesarias2025decoding}
L.~Garc{\'e}s~Arias, S.~Mille, and L.~Padro.
\newblock On the impact of decoding strategies on the detectability of machine-generated text.
\newblock In \emph{COLING}, 2025.

\bibitem[Gerstgrasser et~al.(2024)]{gerstgrasser2024model}
M.~Gerstgrasser, R.~Schaeffer, A.~Dey, R.~Rafailov, H.~Sleiman, J.~Hughes, T.~Korbak, R.~Agarwal, D.~Levine, et~al.
\newblock Is model collapse inevitable? {Breaking} the curse of recursion by accumulating real and synthetic data.
\newblock \emph{arXiv preprint arXiv:2404.01413}, 2024.

\bibitem[Hans et~al.(2024)]{hans2024binoculars}
A.~Hans, A.~Schwarzschild, V.~Cherepanova, H.~Kazemi, A.~Saha, M.~Goldblum, J.~Geiping, and T.~Goldstein.
\newblock Spotting {LLMs} with {Binoculars}: Zero-shot detection of machine-generated text.
\newblock In \emph{ICML}, 2024.

\bibitem[Holtzman et~al.(2020)]{holtzman2020curious}
A.~Holtzman, J.~Buys, L.~Du, M.~Forbes, and Y.~Choi.
\newblock The curious case of neural text degeneration.
\newblock In \emph{ICLR}, 2020.

\bibitem[Kirchenbauer et~al.(2023)]{kirchenbauer2023watermark}
J.~Kirchenbauer, J.~Geiping, Y.~Wen, J.~Katz, I.~Miers, and T.~Goldstein.
\newblock A watermark for large language models.
\newblock In \emph{ICML}, 2023.

\bibitem[Liang et~al.(2024)]{liang2024monitoring}
W.~Liang, M.~Yuksekgonul, Y.~Mao, E.~Wu, and J.~Zou.
\newblock {GPT} detectors are biased against non-native {English} writers.
\newblock In \emph{Patterns}, 4(7), 2024.

\bibitem[Mitchell et~al.(2023)]{mitchell2023detectgpt}
E.~Mitchell, Y.~Lee, A.~Khazatsky, C.~D. Manning, and C.~Finn.
\newblock {DetectGPT}: Zero-shot machine-generated text detection using probability curvature.
\newblock In \emph{ICML}, 2023.

\bibitem[Shumailov et~al.(2024)]{shumailov2024model}
I.~Shumailov, Z.~Shumaylov, Y.~Zhao, N.~Papernot, R.~Anderson, and Y.~Gal.
\newblock {AI} models collapse when trained on recursively generated data.
\newblock \emph{Nature}, 631:755--759, 2024.

\bibitem[Uchendu et~al.(2020)]{uchendu2020authorship}
A.~Uchendu, T.~Le, K.~Shu, and D.~Lee.
\newblock Authorship attribution for neural text generation.
\newblock In \emph{EMNLP}, 2020.

\bibitem[Wang et~al.(2024)]{wang2024editlens}
M.~Wang, Z.~Wang, B.~Li, and J.~Chen.
\newblock {EditLens}: Evaluating and predicting editing on {AI}-generated texts.
\newblock \emph{arXiv preprint arXiv:2404.17435}, 2024.

\bibitem[Wettig et~al.(2024)]{wettig2024qurating}
A.~Wettig, A.~Gao, Z.~Zhong, and D.~Chen.
\newblock {QuRating}: Selecting high-quality data for training language models.
\newblock In \emph{ICML}, 2024.

\bibitem[Xie et~al.(2023)]{xie2023dsir}
S.~M. Xie, S.~Santurkar, T.~Ma, and P.~Liang.
\newblock Data selection for language models via importance resampling.
\newblock In \emph{NeurIPS}, 2023.

\bibitem[Radford et~al.(2019)]{radford2019gpt2}
A.~Radford, J.~Wu, R.~Child, D.~Luan, D.~Amodei, and I.~Sutskever.
\newblock Language models are unsupervised multitask learners.
\newblock \emph{OpenAI Technical Report}, 2019.

\bibitem[Biderman et~al.(2023)]{biderman2023pythia}
S.~Biderman, H.~Schoelkopf, Q.~Anthony, H.~Bradley, K.~O'Brien, E.~Hallahan, M.~A. Khan, S.~Purohit, U.~S. Prashanth, E.~Raff, et~al.
\newblock {Pythia}: A suite for analyzing large language models across training and scaling.
\newblock In \emph{ICML}, 2023.

\bibitem[Groeneveld et~al.(2024)]{groeneveld2024olmo}
D.~Groeneveld, I.~Beltagy, P.~Walsh, A.~Bhagia, R.~Kinney, O.~Tafjord, A.~Jha, H.~Ivison, I.~Magnusson, Y.~Wang, et~al.
\newblock {OLMo}: Accelerating the science of language models.
\newblock In \emph{ACL}, 2024.

\bibitem[Hu et~al.(2022)]{hu2022lora}
E.~J. Hu, Y.~Shen, P.~Wallis, Z.~Allen-Zhu, Y.~Li, S.~Wang, L.~Wang, and W.~Chen.
\newblock {LoRA}: Low-rank adaptation of large language models.
\newblock In \emph{ICLR}, 2022.

\bibitem[Breiman(2001)]{breiman2001randomforests}
L.~Breiman.
\newblock Random forests.
\newblock \emph{Machine Learning}, 45(1):5--32, 2001.

\bibitem[Yang et~al.(2024)]{yang2024qwen2}
A.~Yang, B.~Yang, B.~Hui, B.~Zheng, B.~Yu, C.~Zhou, et~al.
\newblock {Qwen2} technical report.
\newblock \emph{arXiv preprint arXiv:2407.10671}, 2024.

\bibitem[Cao et~al.(2020)]{cao2020coral}
W.~Cao, V.~Mirjalili, and S.~Raschka.
\newblock Rank consistent ordinal regression for neural networks with application to age estimation.
\newblock \emph{Pattern Recognition Letters}, 140:325--331, 2020.

\bibitem[Liang et~al.(2023)]{liang2023helm}
P.~Liang, R.~Bommasani, T.~Lee, D.~Tsipras, D.~Soylu, M.~Yasunaga, Y.~Zhang, D.~Narayanan, et~al.
\newblock Holistic evaluation of language models.
\newblock \emph{TMLR}, 2023.

\bibitem[Hendrycks et~al.(2021)]{hendrycks2021mmlu}
D.~Hendrycks, C.~Burns, S.~Basart, A.~Zou, M.~Mazeika, D.~Song, and J.~Steinhardt.
\newblock Measuring massive multitask language understanding.
\newblock In \emph{ICLR}, 2021.

\bibitem[Achiam et~al.(2023)]{achiam2023gpt4}
J.~Achiam, S.~Adler, S.~Agarwal, L.~Ahmad, I.~Akkaya, F.~L.~Aleman, et~al.
\newblock {GPT-4} technical report.
\newblock \emph{arXiv preprint arXiv:2303.08774}, 2023.

\end{thebibliography}

% ============================================================
% APPENDIX
% ============================================================
\appendix

\section{Wild-Data Calibration}
\label{app:wild}

\begin{table}[h]
\centering
\caption{Depth distribution predictions on wild web-crawl data (10K samples each, Pythia-1.4B reference model). FPR = percentage of samples predicted as non-d0.}
\label{tab:wild}
\small
\begin{tabular}{lcccc|c}
\toprule
\textbf{Corpus} & \textbf{d0 (\%)} & \textbf{d1 (\%)} & \textbf{d2 (\%)} & \textbf{d3 (\%)} & \textbf{FPR (\%)} \\
\midrule
C4         & 78.5 & 19.8 & 1.2 & 0.5 & 21.5 \\
FineWeb    & 82.0 & 16.8 & 0.8 & 0.4 & 18.0 \\
RedPajama  & 85.0 & 13.5 & 1.0 & 0.5 & 15.0 \\
\bottomrule
\end{tabular}
\end{table}

\begin{table}[h]
\centering
\caption{Cross-model agreement on wild-data d0 predictions across three reference models.}
\label{tab:wild_agreement}
\small
\begin{tabular}{lc}
\toprule
\textbf{Metric} & \textbf{Value} \\
\midrule
d0 agreement (3 models) & 92\% \\
Primary misclassification type & template/SEO/boilerplate \\
\bottomrule
\end{tabular}
\end{table}

FPR inversely correlates with corpus cleaning strictness: RedPajama (most filtered) $<$ FineWeb $<$ C4. Misclassified samples are predominantly template-heavy, SEO-optimized, or boilerplate content, whose low lexical diversity and high repetition rates place them in the same feature-space region as shallow recursive generation output.

\section{API Model Depth Estimation}
\label{app:api}

\begin{table}[h]
\centering
\caption{Depth estimation on GPT-4o-mini recursive chains. Fine-tuned: LoRA fine-tuning API chain. Prompted: iterative rewrite prompting without fine-tuning.}
\label{tab:api}
\small
\begin{tabular}{lccc}
\toprule
\textbf{Chain Type} & \textbf{4-class Acc} & \textbf{d0v1 AUC} & \textbf{d1v2 AUC} \\
\midrule
Fine-tuned chain & 58.0 & 0.93 & 0.72 \\
Prompted chain   & 45.0 & 0.78 & 0.61 \\
\midrule
Base model RF (reference) & 78.3 & 0.99 & 0.86 \\
\bottomrule
\end{tabular}
\end{table}

Fine-tuned chain top features: \texttt{p90\_ppl}, \texttt{var\_surprisal}, \texttt{mean\_ppl}, \texttt{p75\_ppl} (3/4 overlap with base models). Prompted chain top features: \texttt{var\_surprisal}, \texttt{p90\_ppl}, \texttt{surprisal\_range}, \texttt{mean\_surprisal} (shifted profile reflecting prompt-mediated constraints). Both settings exceed chance (25\%), confirming that RLHF regularization attenuates but does not eliminate depth signal.

\section{End-to-End vs.\ Feature-Based Methods}
\label{app:dcol}

\begin{table}[h]
\centering
\caption{Comparison of end-to-end neural methods vs.\ feature-based classifiers. DCOL v1--v6 explored contrastive ordinal learning with RoBERTa-large embeddings; all variants underperformed feature-based baselines.}
\label{tab:dcol}
\small
\begin{tabular}{lcc}
\toprule
\textbf{Method} & \textbf{4-class Acc (\%)} & \textbf{Category} \\
\midrule
DeBERTa (fine-tuned)     & 58.3 & End-to-end \\
RoBERTa (fine-tuned)     & 59.5 & End-to-end \\
DCOL v6 (best variant)   & 64.5 & End-to-end \\
\midrule
LR (19 features)         & 76.0 & Feature-based \\
RF (19 features)         & \textbf{78.3} & Feature-based \\
\bottomrule
\end{tabular}
\end{table}

The systematic underperformance of end-to-end methods is attributable to the narrow inter-class boundaries in embedding space. Adjacent recursive depths produce text with cosine similarity exceeding 0.95 in RoBERTa-large embedding space, leaving insufficient margin for contrastive learning to establish discriminative boundaries. Feature-based methods circumvent this by operating on hand-crafted statistics (tail perplexity, surprisal variance) that amplify the distributional differences that embedding-space distances compress. DCOL v1--v6 failure modes included: embedding collapse (v1--v2), out-of-memory (v3), catastrophic forgetting (v4), NaN gradients (v5), and overfitting to depth 0 (v6). This negative result suggests that ordinal depth estimation in text is fundamentally better suited to explicit statistical features than to end-to-end representation learning, at least with current transformer-based text encoders.

\end{document}
